\section{Discussion}
\label{sec:discussion}

\subsection{Held-Out Selection vs.\ Overfitting}
\label{sec:overfit}

The schedule-only interface constrains the LLM to a 4-dimensional
function space---trivial to implement, cheap to evaluate
(30--60\,s CPU)---yet the two winning frontier schedules use
qualitatively different mechanisms (Gaussian bumps vs.\ sinusoidal
shake), suggesting the space is rich.

\emph{Held-out selection materially improves the deployed result.}
For each top run we compare two deployment strategies: deploy the
\textbf{train-selected} script (highest training AEP) vs.\ the
\textbf{held-out-selected} script (highest feasible ROWP AEP among
those the run produced). Across six top runs, held-out selection
gains an average of $+14.4$\,GWh, with gains as large as $+38.2$
(Claude 6hr). For several runs, the training-optimal script is
\emph{infeasible} on ROWP; deploying it without the held-out check
would produce broken layouts (per-run detail:
Appendix~\ref{app:deploy_gap}).

\emph{Overfitting is the default, not the exception.} Of six top
runs, five produced training-optimal schedules that are either
infeasible on ROWP or score below the run's own best-feasible
schedule. LLMs efficiently push training AEP toward the ceiling
and along the way find schedules that exploit the training
polygon's slack. The spacious 5-vertex DEI polygon permits mild
constraint violations without significantly reducing AEP; the
tighter ROWP polygon makes such violations expensive. This
asymmetry---easy training, hard held-out---is not a bug, it is
what makes the held-out signal informative.

\subsection{Why the Held-Out Farm Matters}

Two specific failure modes are invisible without held-out evaluation:

\paragraph{Penalty weight affects quality only weakly.}
A controlled ablation multiplies Claude iter~192's $\alpha$
trajectory by factors from 0.01$\times$ to 10$\times$ while holding
the LR schedule, bumps, and Adam betas fixed. All nine factors
produce feasible ROWP layouts under K=50 stochastic gradient (3/3
init seeds per factor)---contrary to the original hypothesis that
low terminal $\alpha$ causes infeasibility. The effect on
\emph{quality} is small: ROWP AEP feasible-mean ranges from
4262.9\,GWh (factor$\,=\,0.25$) to 4267.9\,GWh
(factor$\,=\,0.01$), a $5.0$\,GWh span across two orders of
magnitude. The schedule's LR dynamics (dual Gaussian bumps,
cosine decay) drive feasibility; the penalty trajectory has only
a weak effect on how efficiently the optimizer uses the available
polygon area.

\paragraph{Training-optimal $\neq$ held-out-optimal.}
Gemini's training-best schedule (AEP = 5600.0) reaches a high
\emph{nominal} ROWP score but produces an infeasible held-out
layout (boundary or spacing constraint violated); after filtering
to feasible layouts only, its best feasible ROWP is 4272.7\,GWh
from a different schedule (training AEP = 5542.2). The schedules
that maximize training AEP do so partly by under-enforcing
constraints---gaining a few GWh from near-boundary placement that
happens to be feasible on DEI but not ROWP. This divergence cannot
be detected without (a)~a physically distinct held-out problem and
(b)~a feasibility filter applied to held-out layouts before
selecting a deployed schedule.

\subsection{Benchmark Properties}

\paragraph{Cost and compute equalization.}
Evaluating a submitted schedule requires only a CPU and takes
30--60 seconds at the training scale ($N\!=\!50$). Running an LLM
agent through the benchmark costs \$5--50 for a 5-hour session
depending on model. The 500-multi-start baseline costs
$\sim$2.5 GPU-hours per farm. Full paper reproduction costs
$\sim$\$500. This is 5--10$\times$ cheaper than SWE-bench
evaluations at comparable scope.

A key design property of the schedule-only interface is
\emph{compute equalization}: every schedule runs the identical
8000-step Adam loop, so evaluation cost depends only on problem
size, not on the schedule's complexity. This eliminates the
confound present in full-optimizer mode, where different strategies
have different compute profiles and the timeout itself shapes which
approaches are viable. In our full-optimizer experiments, 14 of 18
scored scripts exceeded the schedule-only 60\,s budget; without
evaluation feedback, all 19 scripts from a separate run timed out.
The schedule-only interface makes the benchmark \emph{timeout-invariant}:
the same evaluation budget applies to every submission.

\paragraph{Difficulty scaling and cheap-to-expensive transfer.}
The DEI farm variants ($N = 30$--$300$) provide a natural difficulty
axis. Low $N$ is easy (sparse turbines, weak wake interactions);
high $N$ is hard (tight packing, strong constraint pressure,
and increasing evaluation cost---$\sim$30\,s at $N\!=\!50$ vs.\
$\sim$340\,s at $N\!=\!200$). The ROWP held-out farm adds a
cross-physics axis (different turbine, polygon, wind). Crucially,
the LLM's search is conducted entirely in the cheap regime
($N\!=\!50$, 30\,s per evaluation), but the discovered schedules
transfer to the expensive regime ($N\!=\!200$--$300$) without
modification. This \emph{cheap-to-expensive transfer} means the
benchmark's low cost is not a limitation: search happens at
low $N$, deployment happens at high $N$.

\paragraph{Saturated within the schedule-only space.}
Three distinct search strategies (Gemini's hand-found schedule,
Claude's hand-found schedule, and random search at matched compute
inside the LLM-proposed dual-bump family) all reach feasible held-out
scores within a $\sim$4\,GWh band around 4272\,GWh. The two winning
frontier schedules are qualitatively different (dual Gaussian bumps
vs.\ decaying sinusoidal shake), suggesting multiple distinct
mechanism regions exist within this band. We expect that exceeding
this band requires either a richer interface (full optimizer mode
with novel constraint handling) or a different train/test pair.

\subsection{Applicability Beyond Wind Farms}

The schedule-only interface requires three ingredients: (i)~a
differentiable physics simulation, (ii)~constraint functions
expressible as penalties, and (iii)~a gradient optimizer whose
schedule can be separated from its update rule. Any system meeting
these requirements can instantiate the same skeleton: swap the wake
model for a finite-element solver, a molecular dynamics engine, or a
differentiable renderer; keep the 4-scalar control surface
\texttt{(lr, alpha, beta1, beta2)} unchanged. The LLM-discovered
schedule structures---non-monotonic LR dynamics coupled to penalty
ramps---address a meta-problem (how to trade exploration for
constraint satisfaction over a fixed optimization horizon) that is
domain-independent.

Two domains are immediate candidates. \emph{Topology optimization}
(SIMP) uses Adam or MMA with a penalty exponent $p$ ramped from 1 to
3 via a ``continuation strategy'' that is directly analogous to the
$\alpha$ schedule; the 30-year debate over monotonic vs.\ adaptive
continuation~\cite{rojas2015continuation} parallels our finding that
non-monotonic penalty dynamics improve held-out quality.
\emph{Trajectory optimization} for robotics uses differentiable
physics with collision-avoidance penalties and hand-tuned schedule
decay. Both share the same 4-scalar control surface and could be
evaluated with the same train/held-out protocol. Cross-domain
transfer---whether schedules discovered on one physics simulation
generalize to another---is the most important open question this
benchmark raises.

\subsection{Limitations}

\paragraph{Single training farm.}
The LLM trains on one farm (DEI) and generalizes to variants of it
plus one distinct held-out farm (ROWP). The DEI turbine-count sweep
($N=30$--$80$, Sec.~\ref{sec:generalization}) tests only \emph{scaling}
generalization---the polygon, turbine model, and wind resource are all
held fixed---so it cannot distinguish schedule structures that
generalize across physics from structures that happen to work on this
particular DEI--ROWP pair. A stronger test would use multiple training
farms and multiple distinct held-out farms drawn from different wind
regimes, rotor classes, and polygon geometries. We flag this as the
most important next-step extension of the benchmark; the current
train/test pair is best read as an existence proof of the held-out
signal, not as a full characterization of cross-physics
generalization.

\paragraph{Schedule-only mode fixes initialization.}
The skeleton uses wind-aware grid initialization. In full-optimizer
mode, agents improve initialization (hexagonal lattice, K-means
clustering). The schedule-only interface cannot express initialization
innovations, which limits what it measures.

\paragraph{Random-search ablation.}
To isolate the LLM's contribution from the richness of the
parameterized schedule family, we run a random search over a
parameterized schedule family (decay type, perturbation type,
warmup, penalty ramp, Adam betas) at $N=320$ samples, matching
Claude's schedule-only attempt count (Sec.~\ref{sec:ablations}).
This ablation helps distinguish "the LLM found a good point in an
easy space" from "the LLM found a qualitatively different structure
that random search does not reach".

\paragraph{Single-cell training as a design choice.}
The single-cell training protocol surfaces a tension between
training generality and evaluation cleanness. Training on multiple
cells (e.g., varying $N$ or wind rose) could produce more robust
schedules but would weaken the held-out signal: positive
generalization from a diverse training set is expected, while
positive generalization from a single training instance is
noteworthy. We chose the latter to make the generalization claim
falsifiable on every held-out cell individually.

\paragraph{Open-source model gap.}
DeepSeek R1 Distill 32B, the strongest open-source model we
tested, consistently underperforms the 500-multi-start baseline
despite using the same portable scaffold and iterative seed
tournament. Its best schedule uses a structurally simpler
constant-then-cosine decay without the non-monotonic dynamics
(bumps, restarts) that frontier models discover. Running the
same scaffold with two versions of the memory template produced
identical outcomes, confirming the bottleneck is model capability.

\paragraph{Single-agent-class study.}
This paper reports two frontier-tier LLM agents (Claude Code
and Gemini CLI) under their rich CLI scaffolds and one
open-source model (DeepSeek R1 32B) under a portable scaffold. We
attempted additional open-source models (Llama 70B, Qwen 32B)
through the same VLLMRunner scaffold but the
cross-scaffold comparison would have conflated model capability
with scaffold capability, and the open-source runs hit engineering
issues (vLLM context-window limits and dependency incompatibility)
that we did not have time to resolve before the deadline. A
clean multi-tier comparison requires running all models inside the
same scaffold, which we leave to future work.
